\section{Conclusion}
\label{sec:conclusion}

We asked how unstructured, egocentric site evidence can be grounded to a unique element in an
allocentric BIM model, in a cold-start regime with no real labels. The answer is a neuro-symbolic
decomposition organized around an \textbf{ontology-grounded, topology-derived spatial address}. The
oracle analysis proves the architecture sound --- the symbolic layer retains the ground truth in 100\%
of cases and, supplied with the type-conditional address, reorders the pool from Top-1 4.9\% to
78.5\% --- and isolates extraction, not retrieval design, as the bottleneck. The realization study
shows how much of that ceiling survives a real detector and how to consume a noisy address: as a
calibrated soft prior inside a recall-fixed pool (union, never intersection), with one deterministic
specialist lifting the addressable subset from 6.6\% to 67.6\% and selective prediction raising the
answered subset to 80.6\% --- the system knows when to abstain. The depth analysis supplies the
architectural principle both rely on: realizable relational discrimination saturates at one hop, so
depth is compiled into the node and learning is placed at the neural$\to$symbolic interface.

Three findings generalize beyond this system. An address field is deployable only if it is defined in
a frame the evidence carries --- image-recoverability is a design constraint, not an afterthought.
When a mature signal combines with an immature one, union preserves recall and intersection destroys
it. And calibration's payoff here is not reweighting (a measured no-op) but \emph{selective
prediction}: the honest deliverable of a grounding system with downstream consequence is a
coverage--accuracy curve and an explicit defer route. Within the regime we can measure, the structured
spatial address is what makes site-to-BIM grounding realizable, auditable, and transferable; validating
it on real project traces is the immediate next step.
